\centerline{\bf \Huge Комплы. Практика}

\begin{flushright}
        \scriptsize{Эпизод 8. Аска прибывает в Японию}
\end{flushright}

\problem{1}
В результате поворота на $90^{\circ}$ вокруг точки $O$ отрезок $A B$ перешел в $A_1 B_1$. Докажите, что медиана $O M$ треугольника $O A B_1$ перпендикулярна прямой $A_1 B$.

\problem{2} 
$\bigl[\textbf{Отбор Эйлера}\bigr]$
Внутри трапеции $A B C D(B C \| A D)$, где $A D=2 B C$, взята точка $F$, для которой $A B=F B$. Точка $M$ - середина отрезка $F D$. Докажите, что $C M \perp F A$.

\problem{3}
$\bigl[\textbf{Регион Эйлера}\bigr]$
Точка $N$ - середина стороны $B C$ треугольника $A B C$, в котором $\angle A C B=60^{\circ}$. Точка $M$ на стороне $A C$ такова, что $A M=B N$. Точка $K$ - середина отрезка $BM$. Докажите, что $A K=K C$.

\problem{4}
$\bigl[\textbf{ММО, 10.3}\bigr]$
Точка $O$ - центр описанной окружности треугольника $A B C, A H$ - его высота. Точка $P$ - основание перпендикуляра, опущенного из $A$ на прямую $C O$. Докажите, что прямая $H P$ проходит через середину отрезка $AB$.

\problem{5}
$\bigl[\textbf{Сложный тургор, младшие, 4 задача}\bigr]$
На сторонах $B C$ и $C D$ ромба $A B C D$ отмечены точки $P$ и $Q$ соответственно так, что $B P=C Q$. Докажите, что точка пересечения медиан треугольника $A P Q$ лежит на диагонали $B D$ ромба.

\problem{6}
$\bigl[\textbf{ВсОШ, 10.2, 2022}\bigr]$
На стороне $B C$ остроугольного треугольника $A B C$ отмечены точки $D$ и $E$ так, что $B D=C E$. На дуге $D E$ описанной окружности треугольника $A D E$, не содержащей точку $A$, нашлись такие точки $P$ и $Q$, что $A B=P C$ и $A C=B Q$. Докажите, что $A P=A Q$.
